\DocumentMetadata{
  pdfversion=2.0,
  pdfstandard=ua-2,
  lang=en-US,
  tagging=on,
  tagging-setup={math/setup=mathml-SE,tabsorder=structure}
}
\documentclass[11pt,letterpaper]{article}
\usepackage[margin=0.68in,includefoot]{geometry}
\usepackage{newcomputermodern}
\usepackage{amsmath,amscd,mathtools}
\usepackage{tikz,adjustbox}
\providecommand{\square}{\mdlgwhtsquare}
\usepackage[hidelinks]{hyperref}
\hypersetup{
  pdftitle={Topology First-Year Exam, January 2019, Part 1},
  pdfauthor={University of Florida Department of Mathematics},
  pdfsubject={Graduate mathematics examination},
  pdfdisplaydoctitle=true,
  bookmarksopen=true
}
\setcounter{secnumdepth}{0}
\setlength{\parindent}{0pt}
\setlength{\parskip}{0.28em}
\setlength{\emergencystretch}{3em}
\setlength{\leftmargini}{2.15em}
\setlength{\leftmarginii}{2.65em}
\setlength{\leftmarginiii}{3em}
\renewcommand{\labelenumi}{\textbf{\theenumi.}}
\pagestyle{empty}
\raggedbottom
\allowdisplaybreaks
\newenvironment{examproblems}[1]{%
  \begin{enumerate}%
  \setcounter{enumi}{#1}%
  \setlength{\topsep}{0.35em}%
  \setlength{\partopsep}{0pt}%
  \setlength{\itemsep}{0.48em}%
  \setlength{\parsep}{0.18em}%
}{\end{enumerate}}
\newenvironment{examsubparts}{%
  \begin{enumerate}%
  \setlength{\topsep}{0.18em}%
  \setlength{\partopsep}{0pt}%
  \setlength{\itemsep}{0.16em}%
  \setlength{\parsep}{0.1em}%
}{\end{enumerate}}
\newenvironment{examinstructions}{%
  \par\begingroup\itshape
}{%
  \par\endgroup\vspace{0.18em}
}
\makeatletter
\renewcommand\section{\@startsection{section}{1}{0pt}%
  {-1.25ex plus -.25ex minus -.1ex}%
  {0.55ex plus .1ex}%
  {\normalfont\large\bfseries}}
\renewcommand{\@maketitle}{%
  \newpage\null\vskip -1em%
  {\footnotesize\bfseries
    \href{https://www.ufl.edu/}{University of Florida}, \href{https://math.ufl.edu/}{Department of Mathematics}\par}%
  \vskip -0.5em%
  \tagstructbegin{tag=Title}%
    {\LARGE\bfseries\href{https://gma.math.ufl.edu/exams/topology-first-year/}{Topology First-Year Exam}, January 2019, Part 1\par}%
  \tagstructend%
  \vskip 0.25em%
  \tagmcbegin{artifact}%
    \hrule height 1pt%
  \tagmcend%
  \vskip 1em%
}
\makeatother
\title{Topology First-Year Exam, January 2019, Part 1}
\author{}
\date{}
\begin{document}
\maketitle
\thispagestyle{empty}
\begin{examinstructions}
Work all problems.
\end{examinstructions}
\section{Theorems and Proofs.}
\begin{examinstructions}
In this section state and prove each theorem. The proofs should be complete without being tedious.
\end{examinstructions}
\begin{examproblems}{0}
\item
Show that a compact Hausdorff space~\(X\) is normal.
\item
State and prove the Contraction Mapping Theorem.
\item
Give a proof that if~\(X\) is any set, there is no function \(f:X\to\mathscr{P}(X)\) which is onto. Here \(\mathscr{P}(X)\) is the power set of~\(X\) or the set of all subsets of~\(X\).
\item
Suppose that~\(X\) is a Hausdorff space. Show that if \(A\subset X\) is compact, then it is closed.
\item
Suppose that~\(X\) and~\(Y\) are metric spaces and that~\(X\) is compact. Show that if \(f:X\to Y\) is continuous, then~\(f\) is uniformly continuous.
\end{examproblems}
\section{Examples and Minor Proofs.}
\begin{examinstructions}
In this section, give brief examples, counterexamples, or quick proofs.
\end{examinstructions}
\begin{examproblems}{5}
\item
State Urysohn’s Lemma..
\item
Is it possible for there to be a continuous function \(f:[0,1]\to[0,1)\) which is onto?
\item
State the Bolzano–Weierstrass Theorem.
\item
Define a continuous function \(f:C\to[0,1]\) which is onto where~\(C\) is the Cantor set and \([0,1]\) is the unit interval.
\item
Give an example of a space~\(X\) which is \(T_1\) but not Hausdorff.
\item
State the Tietze Extension Theorem.
\item
Define a quotient map.
\item
State the Baire Category Theorem.
\item
Let \(\{X_\alpha\}_{\alpha\in A}\) be a collection of topological spaces. What is the product topology on \(\prod_{\alpha\in A}X_\alpha\)?
\item
What does it mean for a space~\(X\) to be first countable? Show that every metric space is first-countable.
\end{examproblems}
\end{document}
